\section{Notation and preliminaries}\label{sec:preliminaries}

Throughout the paper, $K$ denotes a field of characteristic zero. Polynomial maps and tuples of variables are written as column vectors. For $H\in K[x_1,\ldots,x_n]^n$, we write
\[
 JH=\left(\frac{\partial H_i}{\partial x_j}\right)_{1\leq i,j\leq n}.
\]
Matrix ranks are computed over the corresponding rational function field. We write $\deg_{x_i}f$ for the degree in $x_i$ and $\deg f$ for total degree. The components of $H$ are \emph{linearly independent} if no nonzero constant linear form annihilates $H$. A polynomial map $H$ is called \emph{normalized} if $H(0)=0$. For a map over a polynomial coefficient ring, normalization in specified polynomial variables means that the constant term in those variables is zero in the coefficient ring.

\subsection{Polynomial automorphisms}
Let $R$ be a commutative ring with identity. We denote by $\GA_n(R)$ the group of polynomial automorphisms of $R[x_1,\ldots,x_n]$, represented as polynomial maps and composed by substitution. An elementary automorphism has the form
\[
 e_i(f)=(x_1,\ldots,x_{i-1},x_i+f,x_{i+1},\ldots,x_n),
 \qquad f\in R[x_1,\ldots,\widehat{x_i},\ldots,x_n].
\]
The subgroup generated by these maps is $\EA_n(R)$. Constant translations are products of elementary automorphisms.

\begin{definition}\label{def:tame}
The tame subgroup is
\[
 \TA_n(R)=\langle\GL_n(R),\EA_n(R)\rangle\subseteq\GA_n(R).
\]
An automorphism $F\in\GA_n(R)$ is \emph{tame over $R$} if $F\in\TA_n(R)$. It is \emph{stably tame over $R$} if, for some $m\geq0$,
\[
 F^{[m]}=(F_1,\ldots,F_n,x_{n+1},\ldots,x_{n+m})\in\TA_{n+m}(R).
\]
\end{definition}

Equivalently, the tame group is generated by affine and elementary automorphisms. In particular, a triangular automorphism with unit diagonal coefficients is tame. When $R$ is itself a polynomial ring, tameness over $R$ treats the variables of the coefficient ring as coefficients. We use this distinction in Section~\ref{sec:blocks}.

\subsection{Nilpotence and linear conjugation}
For a $3\times3$ matrix $N$ over a characteristic-zero domain, let $\sigma_2(N)$ denote the sum of its principal minors of order two. The characteristic polynomial yields
\begin{equation}\label{eq:nilpotence}
 N\text{ is nilpotent}
 \quad\Longleftrightarrow\quad
 \tr N=\sigma_2(N)=\det N=0.
\end{equation}
In this case $N^3=0$. The stronger identity $N^2=0$ is not assumed. If $JH$ is nilpotent, then $\det(I+sJH)=1$ for an indeterminate $s$.

For $T\in\GL_n(K)$, set $H^T=T\circ H\circ T^{-1}$. Then
\begin{equation}\label{eq:constant-conjugacy}
 JH^T(x)=T(JH)(T^{-1}x)T^{-1}.
\end{equation}
Thus nilpotence, rank, normalization at the origin, and linear independence of the components are preserved. Formula~\eqref{eq:constant-conjugacy} requires a constant matrix. Nonlinear coordinates below are used either as auxiliary polynomial coordinates or as explicit factors of $\id+sH$.

We use two known triangularization results. A normalized map whose nilpotent Jacobian has rank at most one has linearly dependent components; see~\cite{BCW} or~\cite[Theorems~2.3 and~3.4]{deBondtYan2014}. A homogeneous map in three variables with nilpotent Jacobian is linearly triangularizable~\cite{deBondtEssen}. Consequently, a nonzero homogeneous map $G\in K[x,y,z]^3$ of degree $d\geq1$ with nilpotent Jacobian can be put in the form
\begin{equation}\label{eq:homogeneous-form}
 G=(0,\gamma x^d,B(x,y)),
\end{equation}
where $B$ is homogeneous of degree $d$. For $d=1$, this is ordinary triangularization of a nilpotent matrix.

\subsection{Derivations and slices}
Let $k$ be a field and let $\mathcal A$ be a commutative $k$-algebra. A \emph{$k$-derivation} of $\mathcal A$ is a $k$-linear map $\delta\colon\mathcal A\to\mathcal A$ satisfying the Leibniz rule
\[
 \delta(fg)=\delta(f)g+f\delta(g)\qquad(f,g\in\mathcal A).
\]
Its \emph{ring of constants} is $\ker\delta=\{f\in\mathcal A:\delta(f)=0\}$; this is a field when $\mathcal A$ is a field. An element $s\in\mathcal A$ is a \emph{slice} for $\delta$ if $\delta(s)=1$. A derivation admitting such an element is called a \emph{derivation with a slice}. For example, $\partial_x$ on $k[x,y,z]$ or $k(x,y,z)$ has the slice $x$.
A derivation $\delta$ is \emph{locally nilpotent} if, for every $f\in\mathcal A$, there is an integer $j\geq1$ such that $\delta^j(f)=0$. Having a slice does not by itself imply local nilpotence. 
% In Section~\ref{sec:classification}, the derivation $\Delta$ is defined on a rational function field and has the slice $x$; no local nilpotence of $\Delta$ on that field is assumed.

\subsection{Closed polynomials and relative constants}
Let $k$ be any field.
\begin{definition}\label{def:closed}
A nonconstant polynomial $r\in k[x_1,\ldots,x_n]$ is \emph{closed} if $k(r)$ is relatively algebraically closed in $k(x_1,\ldots,x_n)$.
\end{definition}

In characteristic zero, this is equivalent to $k[r]$ being integrally closed in the ambient polynomial ring~\cite[Lemma~3]{ArzhantsevPetravchuk}. We use the closed-generator theorem: a finite family of polynomials generating a field of transcendence degree one is contained in $k[r]$ for a closed polynomial $r$; see~\cite{ArzhantsevPetravchuk}. Indeed, choose a generative polynomial for one nonconstant member. Every other member belongs to $k(r)$ by relative algebraic closedness, and
\begin{equation}\label{eq:intersection}
 k(r)\cap k[x_1,\ldots,x_n]=k[r].
\end{equation}
For completeness, write a rational function as $A(r)/B(r)$ with coprime $A,B\in k[t]$. A B\'ezout identity for these polynomials remains valid after substitution. If the quotient is polynomial, $B(r)$ divides $A(r)$ in the ambient polynomial ring, so $B(r)$ is a unit. Hence $B$ is constant. This last argument does not require closedness.

The following lemma collects the field-theoretic facts used below and specifies the constant field in the generic-fiber argument.

\begin{lemma}\label{lem:field-facts}
The following statements hold in characteristic zero.
\begin{enumerate}
\item If $\mathcal K/E$ is finitely generated and $L$ is the relative algebraic closure of $E$ in $\mathcal K$, then $L/E$ is finite and $\mathcal K/L$ is regular. Every finitely generated $L$-subalgebra of $\mathcal K$ is geometrically integral.
\item A derivation of a field extends uniquely to a finite separable extension. The extensions of commuting derivations commute.
\item Suppose $\mathcal K/k(t_1,\ldots,t_m)$ is finite separable, and extend the coordinate derivations to $\mathcal K$. If $\bar{a}\in \mathcal K$ is killed by $\partial_{t_i}$ for $i\in I$, then $\bar{a}$ is algebraic over $k(t_j:j\notin I)$.
\end{enumerate}
\end{lemma}

\begin{proof}
For (i), choose a transcendence basis $t_1,\ldots,t_r$ of $\mathcal K/E$ and put $d=[\mathcal K:E(t_1,\ldots,t_r)]$. Every finite intermediate extension $L_0/E$ contained in $L$ satisfies
\[
 [L_0:E]=[L_0(t_1,\ldots,t_r):E(t_1,\ldots,t_r)]\leq d.
\]
A finite subextension of maximal degree therefore contains every element of $L$, which proves finiteness. If an irreducible polynomial over $L$ factors into monic factors over $\mathcal K$, the coefficients of those factors lie in $\mathcal K$ and are algebraic over $L$. They therefore belong to $L$, contradicting irreducibility. Since finite extensions in characteristic zero are separable and simple, $\mathcal K$ is linearly disjoint from every finite extension of $L$. Hence $\mathcal K\otimes_L\overline L$ is a domain. Tensoring an inclusion of an $L$-subalgebra into $\mathcal K$ preserves injectivity and proves geometric integrality.

For (ii), if $\bar{a}$ has minimal polynomial $P(T)=\sum_j c_jT^j$, an extension of $\delta$ must satisfy
\[
 \delta \bar{a}=-\frac{\sum_j(\delta c_j)\bar{a}^j}{P'(\bar{a})}.
\]
Separability makes the denominator nonzero and gives existence and uniqueness. The commutator of two extended derivations is a derivation that vanishes on the base field and is therefore zero.

For (iii), differentiate the monic minimal polynomial of $\bar{a}$ over $k(t_1,\ldots,t_m)$. If $\partial_{t_i}\bar{a}=0$, the differentiated polynomial has smaller degree and still annihilates $\bar{a}$, so it must be zero. Thus its coefficients are killed by every indicated coordinate derivation. Their common constants in the rational function field are $k(t_j:j\notin I)$, which proves the assertion.
\end{proof}

The base field is relatively algebraically closed in a purely transcendental extension. In particular, the common constants of all extended coordinate derivations on a subfield of $k(x,y,z)$ are $k$ whenever the preceding lemma applies.
